\section{الفصل~\ref{sec:sleep}. النوم}

أوضاع النوم والإعادات وتناقص المشابك مُقاسة بتسجيلات العصبونات والديال المجهري والمجهر الإلكتروني؛ وجدول النوم والأرجحة البطيئة توصفهما نماذج الكتاب، أما غرض النوم ففرضيات في الغالب.

{\footnotesize
\setlength{\tabcolsep}{3pt}\renewcommand{\arraystretch}{1.15}
\begin{longtable}{>{\raggedright}p{0.5cm}>{\raggedright}p{4.0cm}>{\raggedright}p{5.6cm}>{\raggedright}p{2.3cm}>{\raggedright\arraybackslash}p{1.7cm}}
\toprule
م & القضية & التجربة وما قيس & المصدر & الحالة \\
\midrule\endfirsthead
\toprule
م & القضية & التجربة وما قيس & المصدر & الحالة \\
\midrule\endhead
1 & في النوم لا تصمت عصبونات القشرة؛ الذي يتغير هو وضع العمل & تسجيلات طويلة لعصبونات القشرة عند الجرذان: في النوم البطيء يبقى التردد غير صفري، وتتناوب فترات العمل مع فترات الصمت؛ وبعد يقظة طويلة يكون التردد أعلى في كل الحالات، وبعد النوم أدنى & \cite{Vyazovskiy2009} & مُقاس \\
2 & \nt{ACET} مرتفع نهارًا وفي النوم السريع، ومنخفض في النوم البطيء (الجدول~\ref{tab:sleepmodes}) & ديال مجهري في القشرة عند قطط تتحرك بحرية: في اليقظة يكون إطلاق الأسيتيل كولين أعلى بـ$100\text{--}175\%$ منه في النوم البطيء، وفي النوم السريع يقارب مستوى اليقظة الهادئة & \cite{Marrosu1995} & مُقاس \\
3 & \nt{NORA} و\nt{SERO} يهدآن في النوم البطيء ويكادان يصمتان في النوم السريع & تسجيل عصبونات \nt{NORA} منفردة عند الجرذان وعصبونات \nt{SERO} عند القطط بحسب مراحل النوم: التردد أقصى في اليقظة، وأدنى في النوم البطيء، ويكاد يكون صفرًا في النوم السريع & \cite{AstonJonesBloom1981,McGintyHarper1976} & مُقاس \\
4 & في النوم السريع تُطفأ العضلات بالتثبيط عبر \nt{GABA} و\nt{GLYC} & عند الجرذان قيس نشاط عصبونات \nt{ACET} لعضلة المضغ وأُعطيت الحاصرات مباشرة إليها: لا يزول الشلل إلا إذا حُجبت مستقبلات \nt{GABA} بنوعيها ومستقبلات \nt{GLYC} معًا & \cite{BrooksPeever2012} & مُقاس \\
5 & ضغط النوم $S$ مكثف بعتبتين~\eqref{eq:sleeppressure}، $\tau_r=18.2$ ساعة، $\tau_d=4.2$ ساعة & نموذج العمليتين؛ الثوابت الزمنية مستخرجة من كيفية ازدياد قدرة الموجات البطيئة في تخطيط الدماغ مع اليقظة وتناقصها في النوم، وشكل العتبات من أوقات الاستيقاظ التلقائي & \cite{Borbely1982,Daan1984} & نظرية \\
6 & تصل الآلة بنفسها إلى $16.1$ ساعة يقظة و$8.0$ ساعات نوم (الشكل~\ref{fig:twoprocess}) & حساب وفق~\eqref{eq:sleeppressure} بعتبات متأرجحة، البرنامج \texttt{models/sleep\_models.py}: $16.1$ ساعة يقظة و$7.9\text{--}8.0$ ساعات نوم & --- & نموذج \\
7 & بعد $40$ ساعة بلا نوم $S=0.90$، لكن النوم لا يدوم إلا $8.8$ ساعة: الدَّين يُسدَّد بالعمق لا بالطول & النموذج: \texttt{models/sleep\_models.py} يعطي $S=0.90$ و$8.8$ ساعة. والتجربة: بعد $40.5$ ساعة بلا نوم عند ثمانية أشخاص كانت قدرة الموجات البطيئة في تخطيط الدماغ في ليلة التعافي الأولى أعلى من المعتاد بدلالة إحصائية & \cite{Borbely1981} & نموذج \\
8 & $S$ فيزيائيًا هو الأدينوزين & ديال مجهري عند القطط: يرتفع الأدينوزين خارج الخلايا في إحدى المناطق المتحكمة في اليقظة خلال اليقظة، ويواصل الارتفاع مع الحرمان من النوم، وينخفض أثناء النوم. أما أن $S$ هو الأدينوزين وحده فلم يُبيَّن & \cite{PorkkaHeiskanen1997,Dunwiddie2001} & فرضية \\
9 & في النوم البطيء تتأرجح القشرة: عمل لبضع مئات من $\text{ms}$ ثم صمت، أقل من مرة في الثانية ($<1\,\text{Hz}$، وتحت التخدير غالبًا $0.3\text{--}0.4\,\text{Hz}$) & تسجيل داخل خلوي عند قطط مخدّرة: إزالة استقطاب مدتها $0.8\text{--}1.5$ ثانية وفرط استقطاب طويل، والإيقاع $<1\,\text{Hz}$، وغالبًا $0.3\text{--}0.4\,\text{Hz}$؛ ويُرى التناوب نفسه بين العمل والصمت في النوم الطبيعي (السطر~1) & \cite{Steriade1993,Vyazovskiy2009} & مُقاس \\
10 & الأرجحة مجموعة بتغذية راجعة وتعب~\eqref{eq:updown}؛ عند $b=5$ الدور $1.1$ ثانية (قيمة النموذج) والعمل نحو $35\%$ من الوقت، وعند $b=1$ العمل منتظم (الشكل~\ref{fig:updown}) & حساب، البرنامج \texttt{models/sleep\_models.py}: الدور $1.11$ ثانية، ونسبة وقت العمل $0.35$ (المتوسط على عدد صحيح من الأدوار)؛ وعند $b=1$ نسبة العمل $1.0$. دور النموذج أقصر من المُقاس (السطر~9) & --- & نموذج \\
11 & \nt{ACET} يطفئ الأرجحة وينقل القشرة إلى العمل المنتظم & تحفيز نواة عصبونات \nt{ACET} عند القط حجب الإيقاع البطيء في $79\%$ من عصبونات القشرة ونقلها إلى العمل المنتظم، أساسًا باختفاء فرط الاستقطاب الطويل. والأسيتيل كولين يثبّط تيارات البوتاسيوم البطيئة التي تصنع التعب & \cite{SteriadeAmzica1993,McCormick1992} & مُقاس \\
12 & ومضات الإيقاع $7\text{--}15\,\text{Hz}$ تولّدها حلقة \foreignlanguage{english}{\nt{GLUT}--\nt{GABA}} بين القشرة وعقدة الترحيل & في شرائح عقدة الترحيل البصرية عند النمس تنشأ الومضات من رشقات ارتدادية لخلايا الترحيل ردًا على تثبيط من نواة \nt{GABA} مجاورة تثيرها هي نفسها & \cite{vonKrosigk1993,SteriadeMcCormickSejnowski1993} & مُقاس \\
13 & إعادات اليوم تجري مسرَّعة: في الحُصين أسرع بنحو $20$ مرة، وفي القشرة بـ$6\text{--}7$ مرات & في النوم البطيء تجري التسلسلات المتكررة لعصبونات الحُصين مضغوطة زمنيًا. بعد الجري على مسار تكررت تسلسلات عصبونات المكان بالترتيب نفسه في ومضات قصيرة من $\sim100\,\text{ms}$، مضغوطة نحو $20$ مرة؛ وفي القشرة الضغط $6\text{--}7$ مرات & \cite{Nadasdy1999,LeeWilson2002,Euston2007} & مُقاس \\
14 & تعزيز تزامن الإعادات مع القشرة يحسّن الذاكرة (علاقة سببية) & عند الجرذان بعد التعلم عُزّز ارتباط الإعادات بالموجات البطيئة وومضات الإيقاع في القشرة بتحفيز مقيَّد بالإعادات: في اليوم التالي كانت الذاكرة عالية، وعند الضوابط بمستوى الصدفة & \cite{Maingret2016} & مُقاس \\
15 & غرض الإعادات هو التعلم المخلوط للذاكرة البطيئة & نظرية نظامي التعلم وحسابات على شبكات اصطناعية: التعلم المتتابع يفسد القديم، والمخلوط لا. ولا توجد تجربة مباشرة على الدماغ تبيّن أن الإعادات لازمة لهذا بالذات & \cite{McClelland1995} & فرضية \\
16 & بعد النوم تتناقص المشابك تناسبيًا مع حجمها، والكبيرة تكاد لا تتغير & مجهر إلكتروني ثلاثي الأبعاد لـ$6920$ مشبكًا في قشرة الفأر: مساحة التماس بعد النوم أصغر بـ$\sim18\%$ مما بعد اليقظة، والتناقص متناسب مع الحجم، والمشابك الكبيرة الثابتة لا تتغير & \cite{deVivo2017} & مُقاس \\
17 & المهمة الرئيسية للنوم إعادة تطبيع الأوزان & فرضية الاتزان الداخلي المشبكي؛ والسطران 1 و16 يتفقان معها، لكنهما لا يثبتان أنها الوظيفة الرئيسية & \cite{TononiCirelli2014} & فرضية \\
18 & النوم السريع اختبار لنموذج العالم على منصة & الوضع (الجدول~\ref{tab:sleepmodes}) مُقاس؛ أما أن الشبكة تشغّل نموذج العالم فافتراض نظري بلا تجربة & \cite{Hobson2009} & فرضية \\
19 & غسل الدماغ في النوم: السؤال مفتوح & تجربة: في النوم يكون الحيز بين الخلايا أوسع بـ$60\%$ والتنظيف أسرع. وأخرى بطريقة قياس مختلفة: في النوم وتحت التخدير يكون التنظيف أبطأ بوضوح. النتائج متناقضة & \cite{Xie2013,Miao2024} & فرضية \\
20 & النوم الموضعي: مناطق من قشرة حيوان مستيقظ تسقط لحظات قصيرة في الصمت & عند الجرذان بعد يقظة طويلة تصمت عصبونات منطقة منفردة من القشرة لحظات قصيرة، كما في النوم البطيء، مع موجة بطيئة في تخطيط الدماغ الموضعي؛ وفي هذه اللحظات يخطئ الجرذ في المهمة أكثر & \cite{Vyazovskiy2011} & مُقاس \\
\bottomrule
\end{longtable}}
